\documentclass[12pt]{article}

\setlength{\parindent}{0pt}
\setcounter{secnumdepth}{3}
\usepackage[margin=1in]{geometry}

% packages
\usepackage{amsmath, amsfonts, amsthm, amssymb}
\usepackage{graphicx, float}
\usepackage{mathtools}
\usepackage{titlesec}
\usepackage{interval}
\usepackage{hyperref}
\usepackage{siunitx}
\usepackage{titling}
\usepackage{vwcol}
\usepackage{setspace}
\usepackage{empheq}
\usepackage{cancel}
\usepackage{esdiff}
\usepackage{multicol}
\usepackage{mdframed}
\usepackage{esdiff}
\usepackage{tikzsymbols}
\usepackage{multicol}
\usepackage{tikz}
\usepackage{varwidth}
\usepackage{pgfplots}
\usepackage{fancyhdr}
\usepackage{enumitem}
\usepackage{tcolorbox}

\intervalconfig {
	soft open fences
}

% header
\pagestyle{fancy}
\fancyhf{}
\lhead{Ethan Fan}
\rhead{AP Calculus BC}
\cfoot{\thepage}

% section format
\titleformat{\section}
  {\large \bfseries}
  {}
  {0em}
  {}
\titlespacing*{\section}{0pt}{0.5em}{0.3em}

\titleformat{\subsection}[runin]
  {\bfseries}
  {}
  {0em}
  {}
\titlespacing*{\subsection}{0pt}{0pt}{0em}

\titleformat{\subsubsection}[runin]
  {\itshape}
  {(\roman{subsubsection})}
  {0em}
  {}

% commands
\newcommand{\alignedintertext}[1]{%
  \noalign{%
    \vskip\belowdisplayshortskip
    \vtop{\hsize=\linewidth#1\par
    \expandafter}%
    \expandafter\prevdepth\the\prevdepth
  }%
}

\newcommand*{\paren}[1]{\ensuremath\left(#1\right)}
\newcommand*{\limit}[2][x]{\ensuremath{\displaystyle\lim_{#1\to#2}}}
\newcommand*{\Limit}[3][x]{\ensuremath{\displaystyle\lim_{#1\to#2}\left[#3\right]}}
\newcommand*{\deriv}[1][x]{\ensuremath{\dfrac{\mathrm{d}}{\mathrm{d}#1}}}
\newcommand*{\Deriv}[2][x]{\ensuremath{\dfrac{\mathrm{d}}{\mathrm{d}#1}\left[#2\right]}}
\newcommand{\ddx}{\dfrac{d}{dx}}
\newcommand{\dydx}{\dfrac{dy}{dx}}
\newcommand{\dydt}{\dfrac{dy}{dt}}
\newcommand{\dxdt}{\dfrac{dx}{dt}}
\newcommand*{\iinteg}[2][x]{\ensuremath{\displaystyle\int #2\;\mathrm{d}#1}}
\newcommand*{\dinteg}[4][x]{\ensuremath{\displaystyle\int_{#2}^{#3}#4\;\mathrm{d}#1}}
\newcommand*{\abs}[1]{\ensuremath{\left|#1\right|}}
\newcommand*{\eps}{\varepsilon}
\newcommand*{\real}{\ensuremath\mathbb{R}}
\newcommand*{\integer}{\ensuremath\mathbb{Z}}
\newcommand*{\posint}{\ensuremath\mathbb{Z^+}}
\newcommand*{\cbrt}[1]{\ensuremath{\sqrt[3]{#1}}}
\newcommand*{\lheqzero}{\ensuremath{\underset{\text{L'H}}{\overset{\left[\frac00\right]}{=}}}}
\newcommand*{\lheqinfty}{\ensuremath{\underset{\text{L'H}}{\overset{\left[\frac{\infty}{\infty}\right]}{=}}}}
\newcommand{\tor}{\text{ or }}
\newcommand{\floor}[1]{\lfloor #1 \rfloor}
\newcommand{\at}{\Big|}

\newcommand{\vem}{\vspace{0.5em}}
\newcommand{\example}[1]{\section{Example #1}}
\newcommand{\problem}[1]{\section{Problem #1}}
\newcommand{\subproblem}[1]{\subsection{(#1)} \,}

\DeclareMathOperator{\dne}{DNE}
\DeclareMathOperator{\sgn}{sgn}

\newcommand{\definition}[1]{\begin{tcolorbox}[colback=red!5!white,colframe=red!75!black,parbox=false] #1 \end{tcolorbox}}

\newcommand{\theorem}[2]{\begin{tcolorbox}[title={#1},colback=blue!5!white,colframe=blue!75!black,parbox=false] #2 \end{tcolorbox}}

\allowdisplaybreaks
% \postdisplaypenalty=100000

% document
\begin{document}

\vspace*{0cm}

% opening
\begin{center}
    {\Large \bfseries Problem Set 24} \\[0.5cm]
    {\large Arc Length} \\[0.35cm]
    {\normalsize November 5, 2026}
\end{center}

% problems

\problem{1}

\subproblem{a}
\begin{gather*}
    n=1 \implies l=\sqrt{2}= 1.414\\
    n=2 \implies l = 1.4604\\
    n=4 \implies l =1.474
\end{gather*}

\subproblem{b}
By the triangle inequality,
\[
AB \le AC+BC
\]
Hence, every polygon length is an underestimate compared to $L$.\vem

\subproblem{c}
\begin{gather*}
    \Delta y=f'(x)\Delta x
\end{gather*}

\subproblem{d}
\begin{gather*}
    y=mx+c\implies y'=m\\
    L=\int_a^b\sqrt{1+m^2}\, dx = (b-a)\sqrt{1+m^2}\\
    y=\sqrt{1-x^2}= \boxed{\frac{\pi}{6}}
\end{gather*}

\problem{2}

\subproblem{b}
\begin{gather*}
    x =y^2\implies x'=2y\\
    l=\int_0^1 \sqrt{1+4y^2}\, dy\\
    \sqrt{1+4y^2} = \sec \theta\\
    y= \frac{1}{2}\tan \theta  \implies dy= \frac{1}{2}\sec^2 \theta \, d\theta\\
    l= \frac{1}{2}\int_0^{\arctan 2} \sec^3 \theta \, d\theta\\
    = \frac{1}{4}(\sec \theta \tan \theta + \ln | \sec \theta + \tan \theta|) \at_0^{\arctan 2}\\
    = \frac{1}{4}(\sqrt{5} \cdot 2 + \ln (2+\sqrt{5})\\
    = \boxed{\frac{\sqrt{5}}{2} + \frac{1}{4} \ln (2+\sqrt{5})}
\end{gather*}

\problem{3}

\subproblem{b}
\begin{gather*}
    x=y+y^3\implies x'=1+3y^2\\
    \boxed{l=\int_1^4 \sqrt{1+(1+3y^2)^2}\,dy}
\end{gather*}

\subproblem{c}
$\sqrt{9+66^2}$ gives the straight-line distance between the two points, functioning as the lower bound. $3+66$ gives the simpler, overestimate version of the integral, functioning as the upper bound.

\problem{4}

\subproblem{b}
\begin{gather*}
    x=\frac{1}{3}\sqrt{y}(y-3)= \frac{1}{3}y^\frac{3}{2}-y^\frac{1}{2}\\
    x'= \frac{1}{2}(y^\frac{1}{2}-y^{-\frac{1}{2}})\\
    l= \int_1^9 \sqrt{1+ \paren{\frac{1}{2}(y^\frac{1}{2}-y^{-\frac{1}{2}})}^2} \, dy\\
    = \int_1^9 \sqrt{1+\frac{1}{4}(y-2+\frac{1}{y})} \, dy\\
    =\int_1^9 \frac{1}{2}\sqrt{ y+2+\frac{1}{y} } \, dy\\
    = \frac{1}{2} \int_1^9 \sqrt{y} + \frac{1}{\sqrt{y}} \, dy\\
    = \frac{1}{3}y^\frac{3}{2} + y^\frac{1}{2} \at_1^9\\
    = 9 +3 - \frac{1}{3}-1 \\
    = \boxed{\frac{32}{3}}
\end{gather*}

\subproblem{c}
\begin{gather*}
    y=\ln \sec x\implies y'=  \sec x \tan x\cdot \frac{1}{\sec x}=\tan x\\
    l=\int_0^\frac{\pi}{4} \sqrt{1+\tan^2x}\,dx\\
    = \int_0^\frac{\pi}{4} \sec x\,dx\\
    = \ln |\sec x +\tan x|\at _0^\frac{\pi}{4}\\
    = \boxed{\ln (\sqrt{2}+1)}
\end{gather*}

\subproblem{e}
\begin{gather*}
    1+(A-B)^2= A^2-2AB+b^2+4AB=(A+B)^2\\
    y= \frac{x^2}{8} - \ln x\\
    y' = \frac{1}{4}x -\frac{1}{x}\\
    l=\int_1^e \sqrt{1+\paren{\frac{1}{4}x -\frac{1}{x}}^2 } \, dx\\
    = \int_1^e \frac{1}{4}x+\frac{1}{x}\, dx\\
    = \frac{x^2}{8} + \ln x\at_1^e\\
    = \boxed{\frac{e^2+7}{8}}
\end{gather*}

\problem{5}
\begin{gather*}
    x^\frac{2}{3}+y^\frac{2}{3}=1 \implies y^\frac{2}{3}=1-x^\frac{2}{3}\\
    y=\paren{1-x^\frac{2}{3}}^\frac{3}{2}\\
    y'= \frac{3}{2} (1-x^\frac{2}{3})^\frac{1}{2}(-\frac{2}{3}x^{-\frac{1}{3}})\\\\
    y'= - \frac{(1-x^\frac{2}{3})^\frac{1}{2}}{x^\frac{1}{3}}\\
    y=x \implies x^\frac{2}{3}=\frac{1}{2} \implies x =\frac{1}{2\sqrt{2}}\\
    l= \int_\frac{1}{2\sqrt{2}}^1 \sqrt{1+ \frac{1-x^\frac{2}{3}}{x^\frac{2}{3}}} \, dx\\
    = \int_\frac{1}{2\sqrt{2}}^1 x^{-\frac{1}{3}}\, dx\\
    = \frac{3}{2}x^\frac{2}{3} \at^1_{\frac{1}{2\sqrt{2}}}\\
    = \frac{3}{2}- \frac{3}{4}\\
    = \boxed{\frac{3}{4}}\\
    8\cdot \frac{3}{4} = \boxed{6}
\end{gather*}
Near $(0,1)$, the derivative approaches negative infinity, rendering the problem more complicated to calculate.

\problem{6}

\subproblem{a}
\begin{gather*}
    y=180-\frac{1}{45}x^2 \implies y'= - \frac{2}{45}x\\
    l= \int_0^{90} \sqrt{1+\frac{4}{2025}x^2}\, dx\\
    x=\frac{45}{2} \tan \theta \implies dx= \frac{45}{2}\sec^2\theta \, d \theta\\
    l=\int_0^{\arctan 4} \frac{45}{2}\sec^3 \theta \, d\theta\\
    = \frac{45}{4} (\sec x \tan x + \ln |\sec x +\tan x |) \at_0^{\arctan 4}\\
    = \frac{45}{4}(\sqrt{17} \cdot 4 + \ln (\sqrt{17} +4) )\\
    = \boxed{45\sqrt{17} + \frac{45}{4} \ln (\sqrt{17} + 4) \approx 209.1\text{m}}
\end{gather*}

\subproblem{b}
The values are the absolute upper and lower bounds to length.

\problem{7}

\subproblem{a}
\begin{gather*}
    y= \int_1^x \sqrt{t^3-1} \, dt \implies y'=\sqrt{x^3-1}\\
    l= \int_1^4\sqrt{1+x^3-1} \, dx\\
    = \frac{2}{5}x^\frac{5}{2} \at_1^4\\
    =\boxed{\frac{62}{5}}
\end{gather*}

\subproblem{b}
\begin{gather*}
    \frac{e^t-e^{-t}}{2}=\int_0^t \sqrt{1+f'(x)^2} \, dx\\
    \frac{e^t+e^{-t}}{2}=\sqrt{1+f'(x)^2}\\
    f'(x)^2= \frac{1}{4}(e^{2t}+2+e^{-2t}-1)\\
    f'(x)=\frac{e^x-e^{-x}}{2}\\
    f(x)= \int \frac{e^x-e^{-x}}{2} \, dx\\
    \boxed{f(x)= \frac{e^x+e^{-x}}{2}}
\end{gather*}

\problem{8}

\subproblem{a}
\begin{gather*}
    y=\frac{2}{3}x^\frac{3}{2} \implies y'=\sqrt{x}\\
    l=\int_0^x \sqrt{1+t} \, dt\\
    =\int_{{1}}^{x+1} \sqrt{u} \, du\\
    = \frac{2}{3}u^\frac{3}{2}\at_1^{x+1}\\
    =\boxed{\frac{2}{3}(x+1)^\frac{3}{2}-\frac{2}{3}}
\end{gather*}

\subproblem{b}
\begin{gather*}
    \frac{2}{3}(x+1)^\frac{3}{2}-\frac{2}{3} = \frac{14}{3}\\
    (x+1)^\frac{3}{2}=8 \implies \boxed{x=3}
\end{gather*}

\subproblem{c}
\begin{gather*}
    s=\frac{2}{3}(x+1)^\frac{3}{2}-\frac{2}{3}\\
    s'= \sqrt{x+1} =\sqrt{3+1}=\boxed{2}
\end{gather*}
It is moving 4 units per second along the curve. 

\problem{9}

\subproblem{a}
\begin{gather*}
    y= \frac{x^3}{6} + \frac{1}{2x} \implies y' = \frac{x^2}{2} - \frac{1}{2x^2} \ge 0\\
    A=\int_1^2 \frac{x^3}{6}+\frac{1}{2x}\, dx\\
    = \frac{x^4}{24}+ \frac{1}{2}\ln x \at_1^2\\
    = \boxed{\frac{5}{8}+\frac{1}{2}\ln 2}
\end{gather*}

\subproblem{b}
\begin{gather*}
    l=\int_1^2 \sqrt{1+ \paren{\frac{x^2}{2} - \frac{1}{2x^2}}^2} \, dx\\
    = \int_1^2 \sqrt{\paren{\frac{x^2}{2} + \frac{1}{2x^2}}^2}\, dx\\
    = \int_1^2 \frac{x^2}{2} + \frac{1}{2x^2}\, dx\\
    = \frac{x^3}{6} - \frac{1}{2x} \at_1^2\\
    = \boxed{\frac{17}{12}}
\end{gather*}

\subproblem{c}
\begin{gather*}
    V=\pi \int_1^2 \paren{\frac{x^3}{6}+\frac{1}{2x}}^2 \, dx\\
    = \pi \int_1^2 \frac{x^6}{36} + \frac{x^2}{6} + \frac{1}{4x^2} \, dx\\
    = \pi (\frac{x^7}{252}+\frac{x^3}{18} - \frac{1}{4x})\at_1^2\\
    = \boxed{ \frac{57\pi}{56}}
\end{gather*}

\subproblem{d}
\begin{gather*}
    V=2\pi \int_1^2 x\paren{\frac{x^3}{6}+\frac{1}{2x}} \, dx\\
    = 2\pi \int_1^2\frac{x^4}{6} + \frac{1}{2} \, dx\\
    = 2\pi ( \frac{x^5}{30} + \frac{1}{2}x)\at_1^2\\
    = \boxed{ \frac{46\pi}{15}}
\end{gather*}

\subproblem{e}
The former represents the straight line distance, which is the strict lower bound, while the latter represent the upper bound, simply the addition of $\Delta x$ and $\Delta y$.

\problem{10}

\subproblem{a}
\begin{gather*}
    L=2\cdot n\cdot \frac{1}{n}=\boxed{2}
\end{gather*}

\subproblem{b}
\begin{gather*}
    l = \int_0^\pi \sqrt{1+n^2\cos^2(n^2 x)} \, dx\\
    \ge \int_0^\pi n|\cos (n^2 x)|\, dx=2n
\end{gather*}

\subproblem{c}
\begin{gather*}
    0 \le f'^2-2f'+1\\
    1+f'^2+2f'\le 2f'^2+2\\
    1+f' \le \sqrt{2} \sqrt{1+(f')^2}
\end{gather*}
It is exactly $\sqrt{2}$ when $f'=1$.\vem

\subproblem{d}
\begin{gather*}
    l=\int_a^b \sqrt{1+f'(x)^2} \, dx\\
    g(x)=k f\paren{\frac{x}{k}} \implies g'(x)=f'\paren{\frac{x}{k}}\\
    l= \int_{ka}^{kb} \sqrt{1+f'\paren{\frac{x}{k}}} \, dx\\
    = k\int_a^b \sqrt{1+f'(u)^2} \, du
\end{gather*}











































\end{document}

